% Generated draft: requires mathematical and source review.
\documentclass[11pt]{article}
\usepackage[margin=25mm]{geometry}
\usepackage{fontspec}
\setmainfont{Latin Modern Roman}
\usepackage{amsmath,amssymb,mathtools,unicode-math}
\setmathfont{Latin Modern Math}
\usepackage{enumitem,needspace,xurl,hyperref}
\setlength{\parindent}{0pt}
\setlength{\parskip}{5pt}
\setlength{\emergencystretch}{3em}
\newcommand{\N}{\mathbb{N}}
\newcommand{\Z}{\mathbb{Z}}
\newcommand{\Q}{\mathbb{Q}}
\newcommand{\R}{\mathbb{R}}
\newcommand{\C}{\mathbb{C}}
\newcommand{\F}{\mathbb{F}}
\newcommand{\A}{\mathbb{A}}
\newcommand{\PP}{\mathbb{P}}
\newcommand{\E}{\mathbb{E}}
\newcommand{\eps}{\varepsilon}
\newcommand{\dd}{\,\mathrm{d}}
\begin{document}
\section*{UnsolvedMath ID: 30003979}
\textbf{Source problem code:} catalogue entry
\subsection{Definitions and mathematical statement}
For notation, let \(k\) denote a base field and let \(\PP^n_k\) denote projective \(n\)-space over \(k\). A projective hypersurface is the zero locus \(V(F)\subseteq\PP^n_k\) of a nonzero homogeneous polynomial \(F\in k[x_0,\ldots,x_n]\) of positive degree. These conventions introduce notation only: the supplied record does not specify the base field, its characteristic, the ambient dimensions, or the degrees under consideration.

\textbf{Research question (catalogue entry).} What general geometric or algebraic mechanism accounts for examples of unexpected projective varieties, including unexpected projective hypersurfaces?

\textbf{Unresolved formulation.} The supplied record does not define ``unexpected,'' specify the conditions imposed on the varieties or hypersurfaces, or give the expected dimension or other benchmark against which unexpectedness is measured. It also does not specify what constitutes a general mechanism, whether every example must arise from such a mechanism, or whether a classification is required. Consequently, this is an open-ended research question as supplied; a precise quantified existence or classification statement cannot be extracted without additional definitions and hypotheses. No particular interpolation model, field, or class of examples is imposed here.
\subsection{Short English statement}
The question asks for a general explanation of the origin of unexpected varieties and hypersurfaces in projective space. The supplied record leaves both ``unexpected'' and the required scope of an explanation undefined, so a precise theorem-like target remains unresolved.
\subsection{Sources}
\begin{itemize}
\item Supplied local UnsolvedMath registry record, permanent ID \(30003979\), title ``Origins of Unexpected Projective Varieties,'' source problem code \texttt{catalogue entry}, in the Oberwolfach Reports Open Problems collection. No individual proposer is supplied; the recorded proposed year is 2018.
\item Supplied source citation: \emph{Mini-Workshop: Asymptotic Invariants of Homogeneous Ideals (2019)}, Oberwolfach Reports, pp.~2739--2767, DOI \texttt{10.4171/owr/2018/45}. \url{https://doi.org/10.4171/owr/2018/45}. The background attributes the question to workshop open-problem discussions documented in Oberwolfach Reports.
\item Supplied literature source labeled ``Unexpected hypersurfaces'': \url{https://arxiv.org/abs/1805.10626}.
\item Supplied literature source labeled ``Line arrangements'': \url{https://arxiv.org/abs/1602.02300}.
\item Supplied literature assessment dated \texttt{2026-08-21}: the record labels the problem open and reports that later work provides classifications and infinite families of unexpected hypersurfaces, while the general origin and classification question remains open. This is supplied dated curation derived from the OWR record, not an independently verified literature assessment; its time-sensitive status and cited claims require independent checking.
\item The supplied statement status is \texttt{corrected\_verified}; its accompanying description says that notation, typography, and sentence structure were repaired while preserving mathematical content. This records the supplied editorial assessment and does not independently establish mathematical precision or correctness.
\end{itemize}
\end{document}
